\clearpage
\section{The uncomfortable part}
\label{sec:uncomfortable}

The following is the research ledger, not a collection of release promises. It distinguishes
fixed defects, competing explanations, experimental code and unrun acceptance tests. The
open-PR inventory at the end preserves the publication snapshot; an open PR is not evidence
that all of its changes are absent from the baseline.

\subsection{The multi-application stall: start with the dependency, not its name}

\readcode The October investigation asks why Windows applications can stall or recover while
individual workloads run well. The later mapping experiments use the Enscape arm of the
SPECviewperf workload. This paper discusses the failure and its timings; it publishes no
SPEC scores, viewset results or proprietary workload images. The experimental mapping line
is distinct from the baseline's earlier relay fixes.\src{admission}\src{rmcost}\src{latency}

A useful causal chain has separate stages: Windows queues work; the relay observes demand;
it publishes host work; the GPU processes a covering marker; host software observes that
marker; the interrupt path notifies the guest; the guest scheduler observes progress. A
silence at the end does not locate the missing or late step. The proposed diagnostic contract
keeps these stages separate and explicitly refuses to identify a private host publication
with a Windows packet without evidence.\src{latency}\src{walknarrow}
\begin{figure}[htbp]
\centering
\begin{tikzpicture}[node distance=15mm]
\node[concept] (a) {Kernel copy work};
\node[concept,right=of a] (b) {Walk + host apply};
\node[concept,right=of b] (c) {Completion observed};
\draw[flow] (a) -- node[above,align=center,font=\scriptsize,text width=25mm] {mapping split
  required} (b);
\draw[flow] (b) -- node[above,align=center,font=\scriptsize,text width=25mm] {result releases work}
  (c);
\end{tikzpicture}
\caption{A serial dependency can resemble a missing interrupt; clocks and packet identity decide.}
\label{fig:stall}
\end{figure}


\subsection{Timeline: defects that really existed, and explanations that did not survive}

\begin{enumerate}[itemsep=0.7em]
\item \readcode The October 9 invalidate incident left a guest trigger incomplete:
  run 223 counted 3,138 invalidates and 3,137 clears. The owner ruling records a guest polling
  that trigger indefinitely. Clearing over absent guest leaves became a policy; refused
  host unmap or host invalidate was explicitly not covered by that exception. A later
  reservation-edge defect was Kayfabe's mapping failure, not invalid guest tables.
  These are documented predecessors, not proof of the later Enscape cause.\src{rulings}
\item \readcode The earlier USER-ring completion gate used GET as a witness. The event-loop
  record says hardware could fetch beyond a dependency before the dependent work ran.
  A tagged processed-reference marker replaced that proof. The attractive explanation
  ``GET passed it, so it executed'' was too strong.\src{eventloop}
\item \measured The README records an earlier relay-handoff campaign with 0 stalled pairs
  out of 26 after the fix, against 3 out of 12 before, plus a later Windows mixed-app session
  of more than fifteen minutes. These are the README's scoped reports at \code{89fc8a222};
  the paper does not claim a population-wide stall rate or transfer that result to the
  later sparse experiment.\src{readme}
\item \readcode A later Translated-slot contention defect consumed a wake while another
  thread held the channel slot. The token could return idle with no owner or future pump.
  The merged recovery retains the owed ring; behavioural tests demonstrated cases that
  source-shape guards had missed. Its design record also says a Turing run still had
  Windows reset storms after the narrow wake defect was fixed.\src{owed}
\item \readcode The first sparse-admission implementation cloned and restored the whole
  reservation tree on failure, returned an uncounted hard error for occupied coarse geometry,
  and could resend a failed acknowledgement on later walks. Local undo, explicit decline
  accounting and retry metadata addressed those software mechanisms. Their CPU tests
  were not an Enscape acceptance run.\src{admission}
\item \measured At sparse-admission revision \code{6d95870c7}, runs 7330 and 7350 showed
  reservation churn near the 8,192-block policy boundary. In run 7330, over approximately
  66 seconds, made objects rose by 241,584 and freed objects by 241,583. The record also
  reports a 1.94-second apply with roughly 59,000 reservation verbs. This refutes the
  simple explanation that every long apply must be one slow native call.\src{rmcost}
\item \measured That same record puts run 7330's first event 153 near the Enscape start
  while recent VA heartbeats were idle. Missing sparse coverage and declines coincided
  with onset. This argues against treating earlier reservation-churn latency as a complete
  account of that particular first reset. It does not prove missing coverage caused it.
  \inferred The hypotheses require separate interventions.\src{rmcost}
\item \readcode Later scope tests passed on the original gate while harvesting a published
  marker and writing GET back before the PUT admission gate. They falsify the narrow
  claim that this gate alone prevents completion observation. They do not falsify delayed
  worker scheduling, lock contention or a dependency before publication.\src{walknarrow}
\end{enumerate}

The timeline is deliberately non-monotone: a real lost wake, a throughput defect and a missing
mapping can all occur during one campaign. Finding the next defect does not retract a
previously demonstrated one, and fixing the previous one does not close the new case.

\subsection{What the later walks actually say}

\measured The public split-latency note attributes \textbf{297,495 of 318,747} walks in boot
7321 and \textbf{461,389 of 490,686} in boot 7391 to one privileged Translated copy-channel
token. These are roughly \textbf{93--94\%}, with one count pair per boot, not a sampled estimate
with a confidence interval. The note is a secondary report supplied to the implementation
review; this paper has not audited the raw host traces.\src{latency}

\measured The scope-narrowing note reports a \textbf{2.43-second apply} and a
\textbf{2.529-second completion gap} in boot 7391 at \code{64c50ad4}. With raised diagnostic
TDR thresholds, that Enscape arm reached its time limit without event 153. The note reports
completion intervals following walk ends in boot 7321. These observations support a
\emph{late rather than permanently lost} account for that arm; they are not a timeout-based
fix, a broad timing distribution, or a one-to-one packet-identity proof.\src{walknarrow}

\inferred A mostly serial channel that asks for a mapping split, waits for its result and
then issues the next request can produce one observed completion near each walk end. That is
consistent with the timing report. Identifying the Windows scheduling node with that exact
host channel still requires ticket-level correlation. Correlation of two periodic series is
not the missing identity.\src{latency}\src{walknarrow}

\readcode A split requests a fresh post-admission walk because releasing later work against
an earlier snapshot could expose stale translations. Same-target tickets admitted together
can share a walk; a request arriving during an active walk belongs to a later submission.
A channel that waits for each result before issuing the next request supplies nothing
concurrent to coalesce.\src{latency}

\textbf{Acceptance still absent from the cited record:} a full production-profile,
default-timeout, sparse-twin-enabled Windows boot that passes the failing mixed workload.
The software validation notes explicitly retain pending host acceptance and the default-OFF
sparse policy. A longer diagnostic timeout cannot fill that gap.\src{sparse}\src{budget}
\src{latency}

\subsection{How our explanation could be wrong}

\textbf{First fork: publication delay versus observation delay.} A packet can spend its time
waiting to be published, executing on hardware, or complete while the observing worker is
busy. The native marker has no GPU-write timestamp in the current diagnostic contract;
\code{marker_gpu_ns=None} is intentional. Publication-to-observed-marker therefore combines
execution and observation latency. Calling all of it ``GPU time'' would be false.
\src{latency}

\textbf{Second fork: CPU scope versus shared hardware.} A VM-wide admission gate can delay an
unrelated ring before publication. Narrowing it by trusted host mutation target removes that
particular dependency. It cannot remove shared engine scheduling, a global host RM lock, or
necessary same-target ordering. If the same stalls persist with unrelated holds removed,
that does not show the scope fix was inert; it shows the intervention did not isolate the
remaining bottleneck.\src{walknarrow}

\textbf{Third fork: trigger, victim and common cause.} A kernel copy channel can request most
walks while an application merely suffers from them. Alternatively, the application's
allocation pattern can cause both sparse work and scheduler distress. Per-channel walk
counts alone cannot distinguish those stories. The decisive observation needs request,
walk start, apply end, readiness, consumption, publication and notification in one life and
clock domain.\src{latency}

\textbf{Proposed falsifier, not a run result:} repeat the same workload and fixed image with
native publication/marker/vector histories and guest event tracing. Include a diagnostic
\code{TdrDelay=60} arm only as an observation window, alongside unchanged default thresholds.
If long completion gaps occur without a corresponding pre-publication split hold, reopen
relay/observation and GPU scheduling explanations. If changing sparse coverage alone changes
onset while latency remains unchanged, prioritise the semantic mapping hypothesis. Record
missing coverage and incomplete timestamps rather than dropping those samples.

\subsection{Sparse backing: semantic fidelity has a resource cost}

\readcode A sparse guest page is not the same as an invalid unmapped address. The selective
sparse-twin experiment exists to reproduce that distinction using host reservation objects.
It stays default OFF. A native micro-oracle can show that the sparse state changes a fault
outcome without proving a long-lived Windows desktop remains correct.\src{sparse}\src{sparsepr}

\readcode The newer policy scales finite per-space and per-session object budgets from the
configured guest framebuffer and counts missing coverage separately from cumulative declines.
Its interval tracker uses an explicit conservative upper bound on overflow; it must not turn
unknown missing coverage into a zero gauge. A permanently excluded host-owned window remains
excluded rather than repeatedly retried. These choices prefer faulting to inventing zero-filled
backing or taking ownership of a host placement.\src{budget}

\inferred This is a design tradeoff that deserves an unfriendly name: the fallback can be
closed for host authority while still incompatible with a guest that needs sparse semantics.
A finite cap protects the host; it does not make the guest's request disappear. The acceptance
question is whether supported workloads stay within the defined mode, not whether a counter
has an appealing value.\src{budget}\src{rulings}

\subsection{Soft budgets are not deadlines}

\readcode The sparse layout allowance limits the work that may \emph{start} in one slice.
Already started replacements must finish or repair their original ownership. An individual
native RM call is not software-preemptible. Ordinary map/unmap, invalidation, retirement and
full teardown are not all covered by the optional-layout allowance.\src{rmcost}\src{budget}

\measured A later CPU simulation with 120-ms releases required 240 ms for release and
mandatory repair, even after two budget bypasses were fixed. That result explicitly disproves
a universal 110-ms expectation. The original assertion of at most one replacement overshoot
was false for error cleanup that started additional work without checking the slice.
These are simulated timings at the scope-repair checkpoint, not a hardware guarantee.
\src{walknarrow}

\readcode The latest experimental continuation defers optional geometry until mandatory
split results and wakes are released, rotates among target identities and stops between
native verbs when a new split arrives. Mandatory repair still completes. Under continuous
mandatory work, optional coverage may remain unhonoured; it is counted rather than hidden.
No hard total-apply deadline or Windows cure is established.\src{latency}

\subsection{Mirrors: retiring memory is an ordering protocol}

\readcode A mirrored address space contains host placements, not merely a dictionary of guest
addresses. The retirement PR delays freeing old backing until invalidation acknowledges that
hardware can no longer use the old translation. A failed invalidate retains ownership.
Its original proposal also keeps explicit unresolved scope tests rather than claiming a
per-range gate exists.\src{retirementpr}\src{walkscope}

\inferred A mirror can be correct after every completed refresh and still produce excessive
latency during a refresh. Conversely, narrowing an admission gate can be fast and wrong if it
releases a source whose mapping is still changing. The experimental target-scoping argument
is about the host object actually mutated, not an assumption that distinct guest roots never
share table pages. Unknown/global requests retain conservative holds.\src{walknarrow}

The remaining fork is per-range certification: can unchanged source spans continue through an
own-space refresh while changed spans permit only old, new or fault outcomes? The latest note
marks that as design work. It needs a bounded change history, root-generation handling,
overflow fallback and publication revalidation; a per-space flag alone is not that protocol.
\src{latency}\src{walknarrow}

\subsection{Pointer mode and lazy observation: the proof costs more than the optimisation}

\readcode Pointer mode exists because repeatedly reading command bytes costs CPU time even
when hardware can safely read the guest's own admitted commands. Registration and migration
still require scanning boundaries. A window anchor is a fact about what was observed before
a reply; retrospectively choosing the current PUT would weaken the proof. The lifecycle
record corrects earlier boundary and mode-scope claims in place.\src{lifecycle}

\readcode Lazy observation additionally makes the guest-visible GET lag. The code compensates
with fresh observations on waits, low credits, split segments, full-ring pressure and mode
changes. Its argument is that staleness can only reduce credit or delay a gate, never create
progress. The small-ring bound matters because a large benchmark ring can hide guest
backpressure.\src{lifecycle}\src{lazycode}

\inferred These mechanisms are justified only if the whole lifecycle remains easier to test
than an always-scanning path is to run. A local counter win is insufficient. One should demand
small-ring tests, hostile mutation schedules, migration at packet boundaries and native
completion timing. The combined scratch performance arms do not reveal which half of lazy
observation deserves its permanent complexity.\src{winperf}\src{lifecycle}

\subsection{Relay debt and recorded disagreement}

\readcode The public ruling records an evolution in the Windows mediation decision: an
initial immutable-copy rule was superseded by the owner's clarification that native USER
privilege plus restricted mappings can permit ordinary GP references unchanged. The earlier
design required immutable command copies; the owner later permitted ordinary pointers under
the native authority boundary. This is a recorded correction to the design, not evidence
that host-owned writable memory became acceptable. Linux Passthrough's stricter mapping prohibition
  remains.\src{rulings}
\src{migrationdecision}

\readcode The latest split note explicitly preserves the existing event-driven relay and
introduces no 64-attempt budget. The new mapping work targets a different dependency. A
convenient retry count would not prove that a lock holder hands progress back correctly.
The contention test's finite fair suffix is an assumption about scheduling, not a product
polling interval.\src{latency}\src{owed}

\unverified This paper found no public repository record sufficient to attribute a current
coordinator preference between every historical relay alternative. It therefore invents no
personal dissent. The practical competing proposals that are recorded are native relocation,
restricted mediation and an optional host-cooperation tier; their acceptance obligations are
more useful than an attributed opinion without a source.\src{migrationdecision}

\subsection{Poor shapes we should be willing to replace}

\begin{description}[itemsep=0.6em]
\item[Global publication dependencies.] \inferred Holding every relayed ring while one
  target changes gives one workload an avoidable latency path into others. The scope
  experiment addresses only this part. Owner policy still requires unchanged VAs to remain
  usable and prohibits solving ordering with a GPU pause.\src{walknarrow}\src{rulings}
\item[Work hidden behind a soft bound.] \readcode A ``100 ms slice'' that excludes repair,
  mapping and teardown is easy to overstate. The newer notes name these exclusions. Owner
  policy requires prompt invalidate completion; changing a diagnostic timer is not compliance.
  \src{rmcost}\src{rulings}
\item[Control-plane scaling.] \readcode The open indexing PR describes a partial correction
  to quadratic RM free paths. It does not claim all graph or hostile-flood costs are linear.
  A cap is a finite memory bound, not a useful throughput bound.\src{indexpr}
\item[Classification by convenient identity.] \readcode The PR catalogue includes separate
  privilege-stamp and mirror-classification work. Guest-handle bit patterns are not an
  authority proof; the rule is native privilege and validated ownership. Integration must
  preserve that rule across old branches, not merely fix the newest path.\src{rulings}
\item[Flag accumulation.] \readcode The flag inventory and cooperative-tier plan distinguish
  measurement switches from product modes. A correctness feature that requires an undocumented
  diagnostic flag has not met the stock-path requirement. The sparse switch currently labels
  an experiment, not a supported mode that must be enabled.\src{flags}\src{cooperation}\src{sparse}
\item[Evidence scattered across heads.] \readcode Status pages, exact-SHA traces and open
  stack PRs often refer to different revisions. The original status page warns about its own
  age. A correction in one page cannot silently update a binary or a downstream summary.
  \src{status}\src{handoff}\src{old}
\end{description}

\subsection{Per-die constants: derive first, then admit the coverage gap}

\readcode The public provenance policy permits computed values, pinned open sources,
unprivileged host queries and source-tagged hardware measurements. Closed-binary material is
not a public data source, build dependency or source of interface layouts. The policy itself
is draft wording awaiting review, and it explicitly does not claim a clean-room process.
This paper carries no findings, values or citations derived from closed binaries.
\src{provenance}\src{credits}

\readcode The constants strategy says to derive or compute first, then ask the unprivileged
host, then measure on the exact die and driver when necessary. A family compression needs both
agreement across measured dies and a mechanism. Agreement of adjacent dies is only
``consistent so far''. An unmeasured die or branch gets a defined, tested refusal or fallback,
not a neighbour's unlabelled row.\src{constants}\src{rulings}

\readcode Coverage is not all on the baseline. The strategy identifies published constants
and display-capability consumers on separate experimental heads while the baseline retains
older paths. Some Windows dies use borrowed display data; Blackwell Windows remains a set of
specific failures and proposed field derivations. Hopper and datacentre cases have no hardware
cell in the guide. A table loader existing is not proof that every caller uses it or that the
fallback runs real applications.\src{constants}\src{requirements}

\inferred A demanding external test is to remove every borrowed optional bit for an unmeasured
die and run the exact acceptance workload. If the fallback only boots but cannot complete
ordinary applications, ``clean mode'' was a name rather than a product configuration. This
follows the owner's explicit fallback requirement, not a new promise of support.\src{constants}

\subsection{Security: unprivileged is a capability limit, not containment evidence}

\readcode Kayfabe trusts the host NVIDIA driver, GPU and VMM adapter. It has no independent
security audit or penetration test. The process has no Kayfabe-owned sandbox; deployment must
confine it. GPU USER privilege and per-channel VA spaces are central because the guest does
not own the physical device. The host IOMMU is a separate host protection layer and does not
provide per-VM identities on this shared physical GPU.\src{security}\src{faq}

\measured The physical-operand oracle at \code{8c084ab7}, RTX 4070 and 595.91.07, reports
that USER copy channels refuse physical operands while virtual operands are delivered.
That is strong evidence for that tested boundary. It is not an audit of every engine,
privilege transition, host-driver bug or same-UID cross-client path.\src{physevidence}\src{faq}

\readcode Deployment recommendations include one UID per VM, dropping capabilities and
confining the whole VMM. The public security page records resource gaps from a dated source
review: GPU time, NVENC sessions, host BAR1, RM call rate and logging were not all bounded per
VM. It explicitly says that review has not been repeated and later code may address parts.
This paper does not turn that historical list into a fresh vulnerability census.\src{security}

\subsection{Fault containment, notification noise and fairness}

\readcode GPU-wide non-stall notifications can wake every registered VM. Extra interrupts
and timing information about engine use are accepted tradeoffs; dropping a real wake would
break correctness. A guest checks its own completion value after a spurious notification.
Optional pacing can delay a signal but must not erase the debt.\src{faq}\src{rulings}

\unverified The effect of a guest-triggered fault on other VMs sharing the same GPU remains
unmeasured in the public guide. Isolation of normal addresses does not establish reset or
fault containment. Likewise a finite framebuffer allocation does not establish fair scheduling,
encoder access or host-call service under an adversarial peer. A useful claim needs a live
victim with retained work across the fault, not just a new process started afterwards.
\src{security}\src{faq}

\subsection{Performance honesty: denominator, placement and missing samples}

\readcode The public Linux performance baseline distinguishes nested rented hosts from the
owner's physical host. It records why the obvious pre-Windows comparison revision could not
run with the current host driver, so the chosen baseline already contains earlier Windows
work. A before/after delta using it cannot measure all Windows cost.\src{linuxperf}

\inferred A throughput ratio is useful only for the named workload and placement. Comparing
a GPU reading host system memory with one reading VRAM answers a placement question as well
as a virtualisation question. An application that batches a token's work cannot grade eager
launch overhead. Three repeated runs in one boot do not estimate boot-to-boot variability.
The older paper's correction record contains exactly these denominator and control failures.
\src{old}\src{linuxperf}

\readcode The display summary warns that average Heaven FPS omits low-frame behaviour and
that its native reference was not repeated with every experimental arm. A fast scratch build
is not a passing integrated lifecycle. No proprietary 3DMark, Steel Nomad or Geekbench score
is reproduced here. No SPEC numerical score is needed to describe the stall.\src{winperf}

\subsection{Unsupported, partially measured and explicitly postponed}

\begin{description}[itemsep=0.6em]
\item[Multiple GPUs.] \readcode A Linux thin guest used two distinct GPUs, sequentially and
  concurrently. CUDA in that guest, UVM across devices, peer access, Windows and choosing a
  display GPU remain outside that result. Calling multi-GPU simply ``unbuilt'' would erase
  real evidence; calling it a complete product would overstate it.\src{status}\src{rulings}
\item[Other VMMs and arm64.] \readcode The owner records these as goals after family coverage.
  QEMU is the implemented adapter. A cross-build or historical pure-core ARM test does not
  establish today's memory/interrupt/reset shell on arm64.\src{rulings}\src{requirements}
\item[Blackwell Windows.] \readcode The guide records GB206 Code 43; open proposals cover
  display fields, constructor admission, native probes and readiness. A successful Linux
  GB205 cell is not a Windows GB206 or GB202 pass.\src{requirements}
\item[Driver unload and recovery.] \readcode The reset and gap documents distinguish device
  life, in-guest reset and reboot. A correct cold boot cannot establish repeated driver
  unload/reload, kill-during-work recovery or generation reuse. Some reset work exists; its
  presence is not a complete lifecycle acceptance.\src{reset}\src{gaps}
\item[Suspend.] \measured Linux S3 passed three cycles at \code{559db30c}; the status record
  reports Windows wake failures. Thus S3 is neither globally absent nor globally working.
  The FSP-family resume path remains a named gap.\src{sleep}\src{handoff}
\item[Display across reboots.] \readcode Reboot scrubbing and boot-layer taint are distinct
  policies: a reboot must scrub before the next OS life, while a driver reset must not
  pretend to create a new trust boundary. The owner's host-allocation isolation premise is
  explicitly accepted, not independently verified by the project.\src{rulings}\src{reset}
\item[Other long tails.] \readcode Multi-head, advanced display fidelity, vIOMMU addressing,
  demand faults, peer resources, container packaging and a rootless launcher need their own
  evidence or implementation. The gap catalogue preserves the issue distinctions; an
  untested configuration is not a supported fallback.\src{gaps}\src{faq}
\end{description}

\subsection{Licensing and provenance are unfinished release work}

\readcode The repository states Apache-2.0 OR GPL-2.0-or-later with exceptions. Some
nvproxy-derived files remain Apache-only; the credits say a QEMU binary linked from the
current tree is not yet distributable under the GPL. The broker has its own notice.
Vendor driver, firmware and OS binaries are not supplied by Kayfabe. These are the project's
stated notices, not this paper's legal interpretation.\src{credits}\src{license}\src{notices}

\readcode Tests obtain vendor software under its own terms. The public downloads note is a
draft and records unresolved installer/automation questions. Building a sealed Windows cache
or automating an installer is engineering work separate from permission to redistribute it.
The policy calls for review of provenance wording and the full history before release.
Nothing in a successful benchmark or build settles those questions.\src{credits}\src{provenance}

\subsection{Process risk: one maintainer, many agents, fallible instruments}

\readcode The repository's workflow assigns an owner approval boundary and uses agents for
research, implementation and review. It requires falsifiers, expected-red tests, multiple
mutation killers, exact revisions and attribution. These mechanisms expose claims to tests;
they do not establish reviewer independence when the same mistaken model informs both code
and oracle. The predecessor records an independent-looking reference tracker with the same
logical omission as the implementation.\src{rules}\src{old}

\inferred The maintainer remains an integration bottleneck: choosing among many experimental
heads requires understanding both the code and which hardware receipts still apply. A
machine-written summary can amplify a confident commit message that was never evidence.
The remedy is source and experiment review, not counting how many reviewers repeat it.
The old paper's corrected object-graph hypothesis illustrates that amplification.\src{old}

\readcode The current PR policy batches real-GPU validation on master and keeps next for
integration. That controls cost, but makes exact final-tree qualification essential. A passing
parent plus a passing child on another date is not a passing merge. Open expected-red PRs and
ignored contracts are useful debt records only while their status remains visible.
\src{rules}\src{owed}

\subsection{Research inventory at the publication cutoff}
\label{sec:inventory}

\readcode The following inventory covers the 86 open PRs captured when writing began, before
this paper's own PR. It includes ordinary and draft PRs, older experiments and a pending batch.
Each entry states its goal, snapshot state and a \emph{proposed} falsifier or acceptance result.
A falsifier here is the paper's research question, not a claim that the PR author preregistered
or ran it. The public PR number and immutable head identify the source. Names from PR titles
are translated into plain descriptions; historical labels are decoded in the appendix.

State ``open'' means GitHub state at capture. ``Head already in baseline'' is a Git ancestry
check, not a claim that a later replacement is identical. Body-reported validation, when
mentioned, has not been rerun. Conflicts, policy decisions and full-profile hardware
qualification remain integration work. The separate experiment notes cited above may contain
later work beyond the older sparse/mirror PR heads.
\paragraph{PR 5. Archive unused predecessor crates without breaking the grader}
\readcode Snapshot: draft, open.
\href{\repobase/pull/5}{Public proposal};
\href{\repobase/commit/c9f00664b439a27b7e52c1157c7262a4c50fa444}{head c9f00664b439}.
\textbf{Proposed falsifier:} The raw-client dependency closure no longer builds or a product crate
  acquires a legacy dependency.

\paragraph{PR 7. Model shared Turing display semaphores}
\readcode Snapshot: draft, open.
\href{\repobase/pull/7}{Public proposal};
\href{\repobase/commit/320b3317033b3755836d59abd69e6b75b2394953}{head 320b3317033b}.
\textbf{Proposed falsifier:} Native acquire/release ordering differs or one window consumes another
  window's event.

\paragraph{PR 8. Compose address-named Blackwell windows inside the store}
\readcode Snapshot: draft, open.
\href{\repobase/pull/8}{Public proposal};
\href{\repobase/commit/598ca60f8b76231e023e2152cc220e07d683c1e3}{head 598ca60f8b76}.
\textbf{Proposed falsifier:} A window escapes the store or a native pixel/ordering reference
  disagrees.

\paragraph{PR 9. Wire exact AD102 Windows initialisation facts}
\readcode Snapshot: draft, open.
\href{\repobase/pull/9}{Public proposal};
\href{\repobase/commit/c3f9c13d557607e40ef95fdaf022268dde875721}{head c3f9c13d5576}.
\textbf{Proposed falsifier:} Native same-die replies or Windows acceptance disagree with the
  supplied fields.

\paragraph{PR 10. Admit measured Turing channel and driver layouts}
\readcode Snapshot: draft, open.
\href{\repobase/pull/10}{Public proposal};
\href{\repobase/commit/728e6febde5e42d53c8085675e2bf1209ab38375}{head 728e6febde5e}.
\textbf{Proposed falsifier:} A same-die driver cell disagrees or an unrelated branch inherits the
  layout.

\paragraph{PR 12. Wire measured GB206 display capabilities across both guests}
\readcode Snapshot: draft, open.
\href{\repobase/pull/12}{Public proposal};
\href{\repobase/commit/24b6a4256ebdd4d806854696d417b93eddf34df4}{head 24b6a4256ebd}.
\textbf{Proposed falsifier:} Exact-head Linux or Windows still fails its application and display
  acceptance.

\paragraph{PR 13. Recover guest command results after agent disconnect}
\readcode Snapshot: draft, open.
\href{\repobase/pull/13}{Public proposal};
\href{\repobase/commit/8aa16cd645572dadeaca065091c6623e7f418540}{head 8aa16cd64557}.
\textbf{Proposed falsifier:} Recovery reexecutes a side-effecting command or associates output with
  another request.

\paragraph{PR 14. Follow early queue movement in the observer}
\readcode Snapshot: draft, open.
\href{\repobase/pull/14}{Public proposal};
\href{\repobase/commit/631f087e34e2e575f303835da9f11fedd3a97df6}{head 631f087e34e2}.
\textbf{Proposed falsifier:} The observer misses a real queue relocation or reads an unvalidated
  candidate.

\paragraph{PR 15. Compare derived display capabilities to native oracles}
\readcode Snapshot: draft, open.
\href{\repobase/pull/15}{Public proposal};
\href{\repobase/commit/886f885c5ffaa9cb49df6058b1b05af61c32371c}{head 886f885c5ffa}.
\textbf{Proposed falsifier:} A mutated wrong capability still passes the claimed behaviour-level
  oracle.

\paragraph{PR 16. Export bounded ordered firmware metadata}
\readcode Snapshot: draft, open.
\href{\repobase/pull/16}{Public proposal};
\href{\repobase/commit/9251e45c45410e62935a01b6e82b0c11a92b58a6}{head 9251e45c4541}.
\textbf{Proposed falsifier:} An overflow looks like complete coverage or ordering changes during
  export.

\paragraph{PR 17. Share CPU views of the backing store}
\readcode Snapshot: draft, open.
\href{\repobase/pull/17}{Public proposal};
\href{\repobase/commit/75b247c07e65cfa48a50141b048170e3827c645f}{head 75b247c07e65}.
\textbf{Proposed falsifier:} View sharing breaks the Linux fast gate or permits stale
  topology/lifetime access.

\paragraph{PR 18. Follow validated public firmware mailbox indirection}
\readcode Snapshot: draft, open.
\href{\repobase/pull/18}{Public proposal};
\href{\repobase/commit/20e0ddc4b0634febb3ed7e97297dc9e31c14705b}{head 20e0ddc4b063}.
\textbf{Proposed falsifier:} An unvalidated pointer becomes a source or a legitimate relocation is
  missed.

\paragraph{PR 19. Require a fresh event for GPU-walk completion}
\readcode Snapshot: draft, open.
\href{\repobase/pull/19}{Public proposal};
\href{\repobase/commit/2bf43d609ff5126d70a1e194c5cb79366a91960a}{head 2bf43d609ff5}.
\textbf{Proposed falsifier:} An old completion can acknowledge a newly submitted walk or a fresh
  result is lost.

\paragraph{PR 20. Stop fast guests without forced process kills}
\readcode Snapshot: draft, open.
\href{\repobase/pull/20}{Public proposal};
\href{\repobase/commit/21b72582ccdbadc771181f6efe0bf05baccffe25}{head 21b72582ccdb}.
\textbf{Proposed falsifier:} The normal stop leaves an owned process alive or kills a process
  outside the run.

\paragraph{PR 21. Retain full startup and teardown trace alignment}
\readcode Snapshot: draft, open.
\href{\repobase/pull/21}{Public proposal};
\href{\repobase/commit/3008bc8346bdfa45162c62c400bffe3f8e7c25e6}{head 3008bc8346bd}.
\textbf{Proposed falsifier:} A truncated trace looks complete or the aligner pairs different device
  lives.

\paragraph{PR 22. Trace replies actually posted to the guest}
\readcode Snapshot: draft, open.
\href{\repobase/pull/22}{Public proposal};
\href{\repobase/commit/4d83e661fcc60f48235381308fd1091dea4351a5}{head 4d83e661fcc6}.
\textbf{Proposed falsifier:} A queued but unposted reply is counted as delivered or a posted reply
  is omitted.

\paragraph{PR 23. Align firmware heartbeat time to the guest timer}
\readcode Snapshot: draft, open.
\href{\repobase/pull/23}{Public proposal};
\href{\repobase/commit/ac01b7bd306dc00202f523cf1545b3cdf9ae026b}{head ac01b7bd306d}.
\textbf{Proposed falsifier:} A changed offset breaks monotonicity or produces a heartbeat in another
  clock domain.

\paragraph{PR 24. Add measured GB205 driver-specific capability words}
\readcode Snapshot: draft, open.
\href{\repobase/pull/24}{Public proposal};
\href{\repobase/commit/5b8dd3a5ff43c73d8434b7cae52e96ab5bc1d9b0}{head 5b8dd3a5ff43}.
\textbf{Proposed falsifier:} Repeat measurement or another branch differs, yet the consumer silently
  reuses the row.

\paragraph{PR 25. Bound metadata traces across guest lifetimes}
\readcode Snapshot: draft, open.
\href{\repobase/pull/25}{Public proposal};
\href{\repobase/commit/fee2ac9336eed69444b9edbb94a6d126d95f80e3}{head fee2ac9336ee}.
\textbf{Proposed falsifier:} Reset evades the cap or old-life records are attributed to a new guest.

\paragraph{PR 26. Select broker guest image and ABI explicitly}
\readcode Snapshot: draft, open.
\href{\repobase/pull/26}{Public proposal};
\href{\repobase/commit/c07f8e8ad4704a0a4060d499493f7ffe3bf3f2e3}{head c07f8e8ad470}.
\textbf{Proposed falsifier:} An implicit image or driver choice overrides the requested test cell.

\paragraph{PR 27. Stop only the owned interactive QEMU safely}
\readcode Snapshot: draft, open.
\href{\repobase/pull/27}{Public proposal};
\href{\repobase/commit/d9aec37d5afe4f208f79945909a5aefc159ca14d}{head d9aec37d5afe}.
\textbf{Proposed falsifier:} Shutdown targets a foreign PID or returns while the owned guest is
  still running.

\paragraph{PR 28. Synchronise bounded guest-agent file staging}
\readcode Snapshot: draft, open.
\href{\repobase/pull/28}{Public proposal};
\href{\repobase/commit/5a300e75d3b1b99fa17b8d4f07a5eea136f4118f}{head 5a300e75d3b1}.
\textbf{Proposed falsifier:} A lost acknowledgement replays a write or a partial transfer is used as
  a complete file.

\paragraph{PR 29. Derive virtual Blackwell tiles and display assignments}
\readcode Snapshot: draft, open.
\href{\repobase/pull/29}{Public proposal};
\href{\repobase/commit/712a0a4755c48fa40eec72724da7d9a09927845f}{head 712a0a4755c4}.
\textbf{Proposed falsifier:} Native acceptance and virtual assignment bounds disagree or an
  unimplemented feature is enabled.

\paragraph{PR 30. Isolate sensitivity to one Blackwell log-buffer refusal}
\readcode Snapshot: draft, open.
\href{\repobase/pull/30}{Public proposal};
\href{\repobase/commit/73af410b683c48b5419651073fb41867f4b21312}{head 73af410b683c}.
\textbf{Proposed falsifier:} Matched refusal and served arms leave the outcome unchanged while
  reachability is proven.

\paragraph{PR 31. Derive exact public display layouts for a driver tag}
\readcode Snapshot: draft, open.
\href{\repobase/pull/31}{Public proposal};
\href{\repobase/commit/98b49a1a4fb6f13bd8bad8214f63039dffa81d42}{head 98b49a1a4fb6}.
\textbf{Proposed falsifier:} The native layout differs or the table is selected for a different
  driver version.

\paragraph{PR 38. Preregister object-binding experiments}
\readcode Snapshot: open, non-draft.
\href{\repobase/pull/38}{Public proposal};
\href{\repobase/commit/ed81e56f37f01f9a78d897b9bf5426d528223d34}{head ed81e56f37f0}.
\textbf{Proposed falsifier:} An arm is interpreted without its positive control or an outcome
  contradicts its recorded prediction.

\paragraph{PR 46. Scope a Windows allocation diagnostic to its measured die}
\readcode Snapshot: draft, open.
\href{\repobase/pull/46}{Public proposal};
\href{\repobase/commit/90364d9d86a634f6319352f42549cacaa77ed9e7}{head 90364d9d86a6}.
\textbf{Proposed falsifier:} An unsupported die or driver inherits the diagnostic answer as a
  product capability.

\paragraph{PR 50. Handle a legacy cursor format and bounded diagnostics}
\readcode Snapshot: open, non-draft.
\href{\repobase/pull/50}{Public proposal};
\href{\repobase/commit/61dc81b73aa2ed25fb07d5b922f88eb450c410e7}{head 61dc81b73aa2}.
\textbf{Proposed falsifier:} Cursor pixels or inversion differ from the native reference, or
  composition violates broker ownership.

\paragraph{PR 51. Derive display DMA bounds from the class}
\readcode Snapshot: draft, open.
\href{\repobase/pull/51}{Public proposal};
\href{\repobase/commit/70afcac481a861def36a4368595524f3a41d8a57}{head 70afcac481a8}.
\textbf{Proposed falsifier:} A boundary operand reaches outside the owned store or a valid last
  element is refused.

\paragraph{PR 52. Observe public Blackwell firmware bootstrap indirection}
\readcode Snapshot: draft, open.
\href{\repobase/pull/52}{Public proposal};
\href{\repobase/commit/b9e2462811ef0310b72fa029c7af2a7325d06289}{head b9e2462811ef}.
\textbf{Proposed falsifier:} A moved queue is missed or an unvalidated location becomes a capture
  source.

\paragraph{PR 53. Admit generated Blackwell video-decode constructors}
\readcode Snapshot: draft, open.
\href{\repobase/pull/53}{Public proposal};
\href{\repobase/commit/aa8fdea9cb0e00fd4e5ec148461d1e805d308e0c}{head aa8fdea9cb0e}.
\textbf{Proposed falsifier:} Decode content or errors differ from the native reference despite
  successful construction.

\paragraph{PR 55. Admit the public Blackwell optical-flow constructor}
\readcode Snapshot: draft, open.
\href{\repobase/pull/55}{Public proposal};
\href{\repobase/commit/844eb167f3a4a5b86e6b001e62ba3f6f7078402b}{head 844eb167f3a4}.
\textbf{Proposed falsifier:} The native constructor or an actual flow workload fails despite an
  accepted class.

\paragraph{PR 56. Review silent registration of a power-management notifier}
\readcode Snapshot: draft, open.
\href{\repobase/pull/56}{Public proposal};
\href{\repobase/commit/c58989dab89e8e033efc9ab807a29771b3d0006c}{head c58989dab89e}.
\textbf{Proposed falsifier:} The guest later requires a hardware action that the accepted inert
  registration did not perform.

\paragraph{PR 57. Add one independently measured GB202 capability scalar}
\readcode Snapshot: draft, open.
\href{\repobase/pull/57}{Public proposal};
\href{\repobase/commit/1bba760c4c6f0d32cb48f9e11a13ccc663786300}{head 1bba760c4c6f}.
\textbf{Proposed falsifier:} A same-die repeat disagrees or the row is applied to an unmeasured
  driver branch.

\paragraph{PR 58. Admit public Blackwell kernel-graphics report releases}
\readcode Snapshot: draft, open.
\href{\repobase/pull/58}{Public proposal};
\href{\repobase/commit/d04f8ae9ad9a6169da7bb2de68bbf0f63d4c6778}{head d04f8ae9ad9a}.
\textbf{Proposed falsifier:} The report value or completion order disagrees with the native oracle.

\paragraph{PR 59. Derive Blackwell post-composition tile capabilities}
\readcode Snapshot: draft, open.
\href{\repobase/pull/59}{Public proposal};
\href{\repobase/commit/680efbb32c7742026812303604cd60e37ede53e1}{head 680efbb32c77}.
\textbf{Proposed falsifier:} A native supported tile is refused or the virtual device advertises a
  tile it cannot execute.

\paragraph{PR 60. Derive Blackwell output-table capacity}
\readcode Snapshot: draft, open.
\href{\repobase/pull/60}{Public proposal};
\href{\repobase/commit/f62ec368f68a204a23b3707a135178cbeb6ba9f9}{head f62ec368f68a}.
\textbf{Proposed falsifier:} The advertised table size exceeds the emulated decoder or rejects a
  valid native boundary.

\paragraph{PR 61. Use allocated family classes in native graphics fixtures}
\readcode Snapshot: draft, open.
\href{\repobase/pull/61}{Public proposal};
\href{\repobase/commit/97dc65cc8ec9a1317bb2dcf9b828c2fb8780688c}{head 97dc65cc8ec9}.
\textbf{Proposed falsifier:} The fixture fails because it uses a foreign class, or false success
  comes from no engine work.

\paragraph{PR 62. Keep process ownership during relay qualification}
\readcode Snapshot: draft, open.
\href{\repobase/pull/62}{Public proposal};
\href{\repobase/commit/fd23ab1c5933d196847eba4c91eae6b97559ed51}{head fd23ab1c5933}.
\textbf{Proposed falsifier:} A harness stop kills a foreign process or loses the GPU lease while its
  owned process is live.

\paragraph{PR 65. Derive post-composition unit widths}
\readcode Snapshot: draft, open.
\href{\repobase/pull/65}{Public proposal};
\href{\repobase/commit/296636c66264b94b591f77fb38dae71d4b13ef1a}{head 296636c66264}.
\textbf{Proposed falsifier:} A boundary-value native transform disagrees or an impossible width is
  advertised.

\paragraph{PR 68. Prototype scoped GPU CI callbacks and endpoint routing}
\readcode Snapshot: draft, open.
\href{\repobase/pull/68}{Public proposal};
\href{\repobase/commit/ffd1d1d518e8550e9b44e070f070ca78e5d37fad}{head ffd1d1d518e8}.
\textbf{Proposed falsifier:} Wrong-run or expired data is admitted, or a URL locator is mistaken for
  content approval.

\paragraph{PR 69. Derive colour-space precision from the open display interface}
\readcode Snapshot: draft, open.
\href{\repobase/pull/69}{Public proposal};
\href{\repobase/commit/6e7cb6ff1befc2c8920a80e615c266901a7044fc}{head 6e7cb6ff1bef}.
\textbf{Proposed falsifier:} A native boundary-value transform disagrees with the derived precision
  or field width.

\paragraph{PR 70. Record unresolved command-descriptor continuation choices}
\readcode Snapshot: draft, open.
\href{\repobase/pull/70}{Public proposal};
\href{\repobase/commit/03715b7d32554995215b477ca24ede10240c5b21}{head 03715b7d3255}.
\textbf{Proposed falsifier:} A packet spanning descriptors is accepted without a consistent
  continuation and lifetime rule.

\paragraph{PR 71. Propagate real translated-channel errors to the guest}
\readcode Snapshot: draft, open.
\href{\repobase/pull/71}{Public proposal};
\href{\repobase/commit/1c6452a276d7014bece3adf42d68cbd4edffbc18}{head 1c6452a276d7}.
\textbf{Proposed falsifier:} A host error becomes success, loses its notifier, or contaminates
  another live channel.

\paragraph{PR 72. Document normal colour decoding and optional diagnostics}
\readcode Snapshot: draft, open.
\href{\repobase/pull/72}{Public proposal};
\href{\repobase/commit/295419c84a542ff2ac1eb428ad5d4ab7cd0800ac}{head 295419c84a54}.
\textbf{Proposed falsifier:} Following the documented default requires a diagnostic flag or produces
  a different transform.

\paragraph{PR 73. Require native privilege stamps for kernel-channel authority}
\readcode Snapshot: draft, open.
\href{\repobase/pull/73}{Public proposal};
\href{\repobase/commit/69ef5973d0440dca5333b13a667557b4c7b1c57b}{head 69ef5973d044}.
\textbf{Proposed falsifier:} Guest-controlled identity bits grant privileged execution without a
  trusted allocation stamp.

\paragraph{PR 74. Split all-space walks at the ABI batch boundary}
\readcode Snapshot: draft, open.
\href{\repobase/pull/74}{Public proposal};
\href{\repobase/commit/e6a0d2b8185dd14e137dd2a06e74d387ee636912}{head e6a0d2b8185d}.
\textbf{Proposed falsifier:} The first space beyond the batch limit is omitted, duplicated or
  prematurely acknowledged.

\paragraph{PR 75. Expose root-move invalidation and withdrawal questions}
\readcode Snapshot: draft, open.
\href{\repobase/pull/75}{Public proposal};
\href{\repobase/commit/01de2fed491aca938da2b862344dafbc56c27d85}{head 01de2fed491a}.
\textbf{Proposed falsifier:} A moved or removed root can acknowledge success while the wrong host
  translation remains live.

\paragraph{PR 76. Admit pointer descriptors using their hardware format}
\readcode Snapshot: draft, open.
\href{\repobase/pull/76}{Public proposal};
\href{\repobase/commit/468a35c9d845afc1938683540acc251bdef90e5a}{head 468a35c9d845}.
\textbf{Proposed falsifier:} Valid hardware descriptor bits are refused or reserved bits create
  extra address authority.

\paragraph{PR 77. Preserve scan fairness and trailing rescan wakes}
\readcode Snapshot: draft, open.
\href{\repobase/pull/77}{Public proposal};
\href{\repobase/commit/1c510cb68e94ee85a717c50c93258bcb6012b1e7}{head 1c510cb68e94}.
\textbf{Proposed falsifier:} One hot token starves another or the final publication is stranded
  without a future wake.

\paragraph{PR 78. Prevent stale ring-tail arithmetic from poisoning publication}
\readcode Snapshot: draft, open.
\href{\repobase/pull/78}{Public proposal};
\href{\repobase/commit/750820e95d88f9f975aa95dbd25edc2e3ed737d7}{head 750820e95d88}.
\textbf{Proposed falsifier:} A wrapped or stale tail underflows into a false capacity or corrupts a
  live producer.

\paragraph{PR 79. Bound streamed command classification on workers}
\readcode Snapshot: draft, open.
\href{\repobase/pull/79}{Public proposal};
\href{\repobase/commit/e018e784cecba94f8a6c43ab0ebc5388271ed5fb}{head e018e784cecb}.
\textbf{Proposed falsifier:} A hostile packet stream consumes unbounded work or a bounded
  continuation loses parser state.

\paragraph{PR 80. Index expensive resource-manager free paths}
\readcode Snapshot: draft, open.
\href{\repobase/pull/80}{Public proposal};
\href{\repobase/commit/5532cdc88ba8b3e4513aafd54ab095aac906b659}{head 5532cdc88ba8}.
\textbf{Proposed falsifier:} A growing object graph still produces quadratic cost in the targeted
  free path.

\paragraph{PR 81. Cap live CPU mapping windows before native placement}
\readcode Snapshot: draft, open.
\href{\repobase/pull/81}{Public proposal};
\href{\repobase/commit/fdccc15e869253f3bfdd1ba0e0da73d150fe2489}{head fdccc15e8692}.
\textbf{Proposed falsifier:} Concurrent admissions exceed the cap or failure rollback leaks
  ownership or a slot.

\paragraph{PR 82. Expose invalidate snapshot drift with failing contracts}
\readcode Snapshot: draft, open.
\href{\repobase/pull/82}{Public proposal};
\href{\repobase/commit/2363313c8c280fbd0d607542f834fa563bc0b16f}{head 2363313c8c28}.
\textbf{Proposed falsifier:} A root change can settle against a stale snapshot; ignored contracts
  are not a runtime fix.

\paragraph{PR 83. Bound guest-driven diagnostics at one admission point}
\readcode Snapshot: draft, open.
\href{\repobase/pull/83}{Public proposal};
\href{\repobase/commit/fa012c5b4d962e841e02b702907d4505409b3a94}{head fa012c5b4d96}.
\textbf{Proposed falsifier:} A flood bypasses the cap through another path or a refusal hides
  unrelated progress.

\paragraph{PR 84. Prevent unsupported device unplug on the current topology}
\readcode Snapshot: draft, open.
\href{\repobase/pull/84}{Public proposal};
\href{\repobase/commit/2165486e0a21985ce204d1776461a88c61af2d51}{head 2165486e0a21}.
\textbf{Proposed falsifier:} The guest can still unplug it or legitimate reset/shutdown loses the
  required lifecycle event.

\paragraph{PR 85. Query display capabilities through unprivileged host requests}
\readcode Snapshot: draft, open.
\href{\repobase/pull/85}{Public proposal};
\href{\repobase/commit/124a37ce3dba62fc70eea47c41bf5a74bf849e1c}{head 124a37ce3dba}.
\textbf{Proposed falsifier:} The query needs privilege or a whole physical answer enables an
  unsupported virtual field.

\paragraph{PR 86. Repair combined-candidate review findings}
\readcode Snapshot: draft, open.
\href{\repobase/pull/86}{Public proposal};
\href{\repobase/commit/ee80c5eb95891874704e33d7f30164bf8c47e26f}{head ee80c5eb9589}.
\textbf{Proposed falsifier:} A hostile-input or ownership regression remains reachable in the exact
  integrated candidate.

\paragraph{PR 87. Check design status and identify dead code}
\readcode Snapshot: draft, open.
\href{\repobase/pull/87}{Public proposal};
\href{\repobase/commit/e50cf4da55d129e985f6c18dc6093d096ada56d1}{head e50cf4da55d1}.
\textbf{Proposed falsifier:} A live proposal is deleted as dead or a source-status correction claims
  implementation exists.

\paragraph{PR 88. Bound trap-side doorbell batches}
\readcode Snapshot: draft, open.
\href{\repobase/pull/88}{Public proposal};
\href{\repobase/commit/c63f66cdbc8633dece1b27d3bfc962f358154fbf}{head c63f66cdbc86}.
\textbf{Proposed falsifier:} An admitted large delta blocks a vCPU or the worker fallback loses or
  duplicates entries.

\paragraph{PR 90. Seed the copy decoder from the allocated host class}
\readcode Snapshot: draft, open.
\href{\repobase/pull/90}{Public proposal};
\href{\repobase/commit/9490d7227e26dc0f5c9c618530ece0c39437cb77}{head 9490d7227e26}.
\textbf{Proposed falsifier:} A family-specific object is decoded with the wrong method set or
  privilege checks weaken.

\paragraph{PR 91. Test translated errors against a live Linux victim}
\readcode Snapshot: draft, open.
\href{\repobase/pull/91}{Public proposal};
\href{\repobase/commit/67df34aff3c7d24e8ba7c2cb9228d63105972b65}{head 67df34aff3c7}.
\textbf{Proposed falsifier:} An induced attacker fault silently succeeds, corrupts the victim, or
  poisons later guest work.

\paragraph{PR 92. Handle mirrored output lookup tables in the colour path}
\readcode Snapshot: draft, open.
\href{\repobase/pull/92}{Public proposal};
\href{\repobase/commit/e636236cf34621b0c9e53170638c54a55a22f979}{head e636236cf346}.
\textbf{Proposed falsifier:} Native colour controls and captured output disagree or the target UI
  operation still freezes.

\paragraph{PR 93. Derive virtual display fields before consulting host facts}
\readcode Snapshot: draft, open.
\href{\repobase/pull/93}{Public proposal};
\href{\repobase/commit/e4cec12b1b31864dfa94b0fca0c1aebbf1726e61}{head e4cec12b1b31}.
\textbf{Proposed falsifier:} A physical host capability enables a virtual feature that the emulated
  display cannot perform.

\paragraph{PR 95. Observe actual Windows adapter readiness independently}
\readcode Snapshot: draft, open.
\href{\repobase/pull/95}{Public proposal};
\href{\repobase/commit/1ac2b51d56ca5fc20bcc08b30c496e20d920635b}{head 1ac2b51d56ca}.
\textbf{Proposed falsifier:} An enable-command return is confused with adapter Code 0 or the
  observer reads a stale instance.

\paragraph{PR 97. Register native completion listeners before submission}
\readcode Snapshot: draft, open.
\href{\repobase/pull/97}{Public proposal};
\href{\repobase/commit/f4879bda249946da0ce237edd84cfce442ed6480}{head f4879bda2499}.
\textbf{Proposed falsifier:} A completion arriving immediately after submission is missed or
  attributed to another client.

\paragraph{PR 98. Ratchet source width at one hundred columns}
\readcode Snapshot: draft, open.
\href{\repobase/pull/98}{Public proposal};
\href{\repobase/commit/820db0c11d16e150d256a522495f895e1da1bba6}{head 820db0c11d16}.
\textbf{Proposed falsifier:} A newly long non-exempt source line passes, or an unchanged allowed
  line creates false debt.

\paragraph{PR 99. Qualify a combined descriptor and mapping candidate}
\readcode Snapshot: draft, open.
\href{\repobase/pull/99}{Public proposal};
\href{\repobase/commit/714ec29d3fdca0a013d9d478755346b0f1af6fa7}{head 714ec29d3fdc}.
\textbf{Proposed falsifier:} The exact combined tree fails full-profile Windows or other-die
  regression, despite parent passes.

\paragraph{PR 101. Inventory stale documentation with a ratchet}
\readcode Snapshot: open, non-draft.
\href{\repobase/pull/101}{Public proposal};
\href{\repobase/commit/b7fd84011a7b2b10cd12a6de7456b432a8bbd2b2}{head b7fd84011a7b}.
\textbf{Proposed falsifier:} A planted stale active document escapes detection or a historical
  archive is treated as live.

\paragraph{PR 103. Correct two historical status checkpoints}
\readcode Snapshot: open, non-draft.
\href{\repobase/pull/103}{Public proposal};
\href{\repobase/commit/19dcdd9b3d50d6453a933921f96f60667002d7f3}{head 19dcdd9b3d50}.
\textbf{Proposed falsifier:} The new status still implies the old checkpoint is the current
  integrated implementation.

\paragraph{PR 107. Calibrate the Windows readiness observer}
\readcode Snapshot: draft, open.
\href{\repobase/pull/107}{Public proposal};
\href{\repobase/commit/95dac5d5eacb3d2f7b9b8763c4f9df556056599f}{head 95dac5d5eacb}.
\textbf{Proposed falsifier:} A slower observer merely hides a product failure or accepts readiness
  without the required state.

\paragraph{PR 108. Measure display fields with a privileged read-only tool}
\readcode Snapshot: draft, open.
\href{\repobase/pull/108}{Public proposal};
\href{\repobase/commit/2ace3d219ad4490af41c5eccff5b48ea7009a228}{head 2ace3d219ad4}.
\textbf{Proposed falsifier:} The tool alters state or its independently decoded replies disagree
  with the native oracle.

\paragraph{PR 110. Retire failed provider machines before retrying}
\readcode Snapshot: draft, open.
\href{\repobase/pull/110}{Public proposal};
\href{\repobase/commit/d53cfd9f44d7bdbc87887c2e19fb3431545d84b5}{head d53cfd9f44d7}.
\textbf{Proposed falsifier:} Retry starts before absence is verified or chooses the same failed
  machine or host.

\paragraph{PR 112. Consume measured Blackwell display capabilities}
\readcode Snapshot: draft, open.
\href{\repobase/pull/112}{Public proposal};
\href{\repobase/commit/008c97705a2c1cd043f407dc1b5c9e8eb446ad1f}{head 008c97705a2c}.
\textbf{Proposed falsifier:} Same-die native answers disagree, or the exact-head Windows A/B still
  fails the target path.

\paragraph{PR 127. Record a Blackwell driver-readiness refusal}
\readcode Snapshot: open, non-draft.
\href{\repobase/pull/127}{Public proposal};
\href{\repobase/commit/f0fa7f4953b9b6c1d651cbf33a8ddef657919ad3}{head f0fa7f4953b9}.
\textbf{Proposed falsifier:} The operator labels a failed provider boot as a tested product cell or
  fails owned cleanup.

\paragraph{PR 137. Prototype a private controller without the Worker}
\readcode Snapshot: draft, open.
\href{\repobase/pull/137}{Public proposal};
\href{\repobase/commit/794d4a700d752018aedf45a28a0df77792effd2a}{head 794d4a700d75}.
\textbf{Proposed falsifier:} A real owned-instance lifecycle violates approval, spend, exact-SHA or
  teardown constraints.

\paragraph{PR 140. Keep retired backing until invalidation acknowledges}
\readcode Snapshot: draft, open.
\href{\repobase/pull/140}{Public proposal};
\href{\repobase/commit/edfd0271f38111b6ba32b14d7fd6cf78d8deaeb9}{head edfd0271f381}.
\textbf{Proposed falsifier:} Old backing is freed while an old translation can execute, or failed
  invalidate loses ownership.

\paragraph{PR 146. Represent guest sparse pages with selective host reservations}
\readcode Snapshot: draft, open.
\href{\repobase/pull/146}{Public proposal};
\href{\repobase/commit/3e52e0161c24547bf2e68d90e79b07c26a30df6d}{head 3e52e0161c24}.
\textbf{Proposed falsifier:} Sparse-ON repeats the failure under the full production profile or
  breaks another die.

\paragraph{PR 160. Bring up the TU106 Windows path}
\readcode Snapshot: draft, open.
\href{\repobase/pull/160}{Public proposal};
\href{\repobase/commit/0b44315ef7f33dd19f919f2937b07c230e73e37a}{head 0b44315ef7f3}.
\textbf{Proposed falsifier:} An exact-head guest fails the full workload despite a successful
  adapter initialisation.

\paragraph{PR 163. Permit an explicit headless Windows device experiment}
\readcode Snapshot: draft, open.
\href{\repobase/pull/163}{Public proposal};
\href{\repobase/commit/e5b9ee2fe647834f9dc141540c64c9792f884b89}{head e5b9ee2fe647}.
\textbf{Proposed falsifier:} Display-disabled setup still creates a display device or is mistaken
  for full desktop acceptance.

\paragraph{PR 164. Model Turing display-window semaphore releases}
\readcode Snapshot: draft, open.
\href{\repobase/pull/164}{Public proposal};
\href{\repobase/commit/c7e085b3d82a6ff01561887c3530eb47d5133bfc}{head c7e085b3d82a}.
\textbf{Proposed falsifier:} The native release oracle and guest notification disagree; acquire
  semantics remain unmodelled.

\paragraph{PR 165. Reduce repeated Windows lease setup}
\readcode Snapshot: draft, open.
\href{\repobase/pull/165}{Public proposal};
\href{\repobase/commit/a1cd0e6283e6261d56164089ce5dac2a6251f4fd}{head a1cd0e6283e6}.
\textbf{Proposed falsifier:} A cached guest differs from the pinned clean image, leaks private
  state, or fails a cold control.

\paragraph{PR 169. Explain a consumed completion wake}
\readcode Snapshot: draft, open.
\href{\repobase/pull/169}{Public proposal};
\href{\repobase/commit/6adec88f9a5f0b847084fbcc8f3d0365b49faa0a}{head 6adec88f9a5f}.
\textbf{Proposed falsifier:} An induced lost wake remains invisible, or a harmless pending channel
  is called stranded.

\paragraph{PR 172. Simplify the GPU trigger workflow}
\readcode Snapshot: open, non-draft.
\href{\repobase/pull/172}{Public proposal};
\href{\repobase/commit/5ced95e19e8dc6df4a996470703e996d642da0d3}{head 5ced95e19e8d}.
\textbf{Proposed falsifier:} Removing the action dependencies changes which approved SHA is tested
  or skips required work.

\paragraph{PR 176. Batch the documentation and site changes}
\readcode Snapshot: open, non-draft; head already in baseline.
\href{\repobase/pull/176}{Public proposal};
\href{\repobase/commit/12c4008705f17033a803e56d4d479bf1275cdb94}{head 12c4008705f1}.
\textbf{Proposed falsifier:} The final batch tree differs from the approved tree or a publication
  check fails.



\subsection{A gap ledger can itself be stale}

\readcode The feature-gap catalogue mixes closed rows, open questions and corrections. Its
priority rule puts application blockers first and console fidelity later. The handoff also
retains superseded decisions, including a capability allowlist already approved and a
migration alternative later tested natively. Keeping history is valuable only if readers
follow dated corrections before reading old verdicts.\src{gaps}\src{handoff}\src{rulings}

\inferred The paper has the same failure mode. Its source pin prevents silently moving the
implementation under a claim; it cannot prevent a new experiment tomorrow from falsifying an
explanation. The inventory should be regenerated for a later edition, and the case study
should change only when a discriminating measurement changes it. The unanswered questions
that follow are intended to help make that happen.
